\NewDocumentCommand{\freigestellterAdlerZeichnen}{O{} r() O{dlrgSchwarz}}{
    \IfFileExists{dlrg_angeschnittener_adler.tex}{
        \input{dlrg_angeschnittener_adler.tex}
        \adlerZeichnen{#1}{#2}{#3}
    }{
        \adlerPlatzhalterZeichnen{#1}{#2}{#3}
    }
}

\NewDocumentCommand{\adlerPlatzhalterZeichnen}{m m m}{
    \begin{scope}[shift={(#2)},y=1pt, x=1pt,yscale=-1,cm={{1,0.0,0.0,-1,(-920,440)}},#1]
        \draw[#3,line width=10pt](500,0) -- (1100,100) -- (1130,140)  -- (900,180) -- (910,240) -- (1230,200) -- (1350,300)-- (1100,800) -- (680,870) -- (790,700) -- (650,800) --(490,750) -- (800,230) -- cycle;
    \end{scope}
}

\NewDocumentCommand{\freigestellterAdlerWasserzeichen}{O{}D(){0cm,8cm}}{
    \begin{tikzpicture}[remember picture,overlay]
        \node[shift={(#2)}] at (current page.south east){
            \begin{tikzpicture}
                \freigestellterAdlerZeichnen[opacity=.25,scale=0.7,#1](0,0)
            \end{tikzpicture}
        };
    \end{tikzpicture}
}
